\chapter{Ответы к задачам}
\label{app:answers}

Здесь собраны ответы и краткие решения к задачам из глав, в порядке книги.

\section*{Сигналы}
\begin{enumerate}
  \item \stars{1}~$101101_2=32+8+4+1=45$; $11111111_2=255$ ($2^8-1$).
  \item \stars{1}~$37=100101_2$ ($32+4+1$); $200=11001000_2$ ($128+64+8$).
  \item \stars{1}~$2^{10}=1024$ значения; для 1000 предметов нужно 10 бит ($2^9=512<1000\le 1024$).
  \item \stars{2}~Уровень не определён: 1,4~В лежит между порогами 0,8 и 2,0~В. Такой сигнал может быть прочитан и как 0, и как 1, поэтому схема ведёт себя непредсказуемо.
  \item \stars{1}~$T = 1/27\,\text{МГц}\approx 37$~нс; за секунду проходит $27\,000\,000$ тактов.
  \item \stars{2}~Сигнал меняется на каждом втором фронте: его частота вдвое ниже, период вдвое больше. Проверка в Logra: урок «Делитель частоты на 2», вкладка Диаграммы.
  \item \stars{2}~Кнопка даёт 1 при нажатии, а светодиод платы горит при 0 (порт \code{led\_n}). Прямое соединение даёт обратное поведение. Исправление: поставить между ними инвертор (элемент НЕ), тогда нажатие зажигает светодиод.
  \item \stars{2}~Ошибка в весах: старший разряд стоит слева, и вес растёт справа налево. Верно: $1\cdot8+0\cdot4+1\cdot2+1\cdot1=11$.
  \item \stars{3}~Схема: door $\to$ НЕ $\to$ \code{led\_n}. Ход решения: светодиод горит от нуля, значит нужна инверсия. Как проверить в Logra: вкладка Диаграммы, \code{led\_n} должен быть точной инверсией door; пример \code{button} (прямое соединение) показывает «неправильный» вариант, пример \code{not} --- правильный.
  \item \stars{3}~Число $37=100101_2$, единицы в разрядах 5, 2 и 0. Светодиоды горят от нуля, поэтому на \code{led\_n} нужен инвертированный код: $011010_2$ (то есть $26$). Как проверить в Logra: шина шести бит, блок НЕ шириной 6; при вводе 37 горят светодиоды разрядов 5, 2 и 0, перебор 0--63 даёт совпадение с таблицей.
\end{enumerate}

\section*{Системы счисления и коды}
\begin{enumerate}
  \item \stars{1}~$\text{B7}_{16}=10110111_2=183$; $\text{3C}_{16}=00111100_2=60$.
  \item \stars{1}~$1998 \to 0001\ 1001\ 1001\ 1000$ (цифры 1, 9, 9, 8 по четыре бита).
  \item \stars{2}~$p_1=0\oplus1\oplus0=1$, $p_2=0\oplus1\oplus0=1$, $p_3=1\oplus1\oplus0=0$; слово (позиции 1--7) \code{1100110}. Проверка: пример \code{hamming-encode}.
  \item \stars{2}~$s_1=1$, $s_2=1$, $s_3=1$, синдром $111_2=7$: ошибка в седьмой позиции. После исправления \code{1100110}, данные $0110$ (позиции 3, 5, 6, 7).
  \item \stars{1}~Четыре единицы, поэтому бит чётности равен 0.
  \item \stars{1}~$0010\,1010 = 2 + \tfrac{10}{16} = 2{,}625$.
  \item \stars{2}~$1000+0101=1101$ --- это 13, но не цифра BCD (больше 9). Прибавляем поправку $0110$: $1101+0110=1\,0011$, то есть перенос 1 и цифра 3: BCD-запись \code{0001\,0011}.
  \item \stars{2}~Бит чётности лишь сообщает, что число единиц нечётно, но не указывает позицию ошибки: неизвестно, какой бит переворачивать. Исправлять умеет код Хэмминга, у которого синдром равен номеру позиции.
  \item \stars{3}~Выход: $d_3d_2d_1d_0p$ с $p=d_3\oplus d_2\oplus d_1\oplus d_0$. Пример: цифра 7 (\code{0111}) даёт $p=1$, слово \code{01111}. Как проверить в Logra: урок «Проверка чётности слова», пример \code{parity}; на вкладке Диаграммы перебрать цифры 0--9, число единиц на выходе должно быть чётно.
  \item \stars{3}~Старшая половина \code{0101}: $p_1=0$, $p_2=1$, $p_3=0$, слово \code{0100101}; младшая \code{1010}: $p_1=1$, $p_2=0$, $p_3=1$, слово \code{1011010}. Испортив позицию 3 в старшем слове, получаем \code{0110101}: $s_1=1$, $s_2=1$, $s_3=0$, синдром $011_2=3$. Переворачиваем бит 3, данные (позиции 3, 5, 6, 7) равны \code{0101}. Как проверить в Logra: примеры \code{hamming-encode} и \code{hamming-syndrome}, вкладка Диаграммы.
\end{enumerate}

\section*{Логические элементы и булева алгебра}
\begin{enumerate}
  \item \stars{1}~$F=1$ для комбинаций $ABC = 001,\ 011,\ 110,\ 111$; для остальных четырёх $F=0$.
  \item \stars{1}~$A$. Это закон поглощения: $A + AB = A(1+B) = A$.
  \item \stars{2}~$\overline{A}\cdot(\overline{B}+\overline{C})$. Ход: $\overline{A+X}=\overline{A}\cdot\overline{X}$, затем $\overline{B\cdot C}=\overline{B}+\overline{C}$.
  \item \stars{2}~НЕ: оба входа вместе (1 элемент). И: И-НЕ и за ним НЕ (2 элемента). ИЛИ: инвертировать оба входа и подать на И-НЕ (3 элемента). Проверка в Logra: урок «И-НЕ из блоков».
  \item \stars{2}~$A\overline{B}+\overline{A}B$. Ход: исключающее ИЛИ равно 1, если входы различны.
  \item \stars{1}~$2^6=64$ строки.
  \item \stars{2}~Ошибка: инверсия суммы не равна сумме инверсий. При $A=1$, $B=0$ левая часть $\overline{1+0}=0$, правая $0+1=1$. Верный закон: $\overline{A+B}=\overline{A}\cdot\overline{B}$.
  \item \stars{2}~При входах 11: ИЛИ даёт 1, а исключающее ИЛИ даёт 0. На остальных комбинациях (00, 01, 10) они совпадают. Проверка в Logra: пример \code{gates}, вкладка Диаграммы.
  \item \stars{3}~Таблица: единица в строках $(\text{door},\text{night},\text{armed})=101,\ 110,\ 111$. Схема: ИЛИ(night, armed), затем И с door, затем НЕ для светодиода платы. Как проверить в Logra: вкладка Диаграммы, перебор восьми комбинаций; светодиод горит ровно в трёх.
  \item \stars{3}~Четыре элемента И-НЕ: $x=\overline{\text{night}}$ и $y=\overline{\text{armed}}$ (два НЕ из И-НЕ), $z=\overline{x\cdot y}=\text{night}+\text{armed}$, $\text{led\_n}=\overline{\text{door}\cdot z}$. Ход: И-НЕ последнего каскада сразу даёт инверсию, нужную светодиоду платы. Как проверить в Logra: диаграмма должна совпасть с диаграммой схемы из задачи~9.
\end{enumerate}

\section*{Комбинационные схемы}
\begin{enumerate}
  \item \stars{1}~$F_1=\overline{B}$ (единицы в строках 0, 1, 4, 5, то есть всюду, где $B=0$); $F_2=C$; $F_3=AB+AC+BC$ (мажоритарная функция). Проверка: соберите в Конструкторе и переберите 8 комбинаций на вкладке Диаграммы.
  \item \stars{1}~$5\times 1{,}5=7{,}5$~нс. Период такта не может быть короче этой задержки (плюс время установки), поэтому частота не выше примерно 130~МГц.
  \item \stars{2}~$F = D\cdot(N+S)$: схема из элемента ИЛИ и элемента И. Ход: условие «и» связывает дверь с суммой «ночь или охрана»; раскрыв скобки, получим $DN+DS$ (три элемента), скобочная форма экономнее.
  \item \stars{1}~Комбинационные: сумматор и мультиплексор. Последовательностные: счётчик и кнопка с переключением, им нужна память.
  \item \stars{2}~Да, при переключении $A$ на выходе возможна короткая «иголка» нуля: $AB$ гаснет, а $\overline{A}C$ ещё не загорелось. Добавляют слагаемое $BC$ (оно держит единицу при $B=C=1$): $F = AB + \overline{A}C + BC$. На карте Карно группы $AB$ и $\overline{A}C$ соприкасаются, но не перекрываются.
  \item \stars{2}~Ответ ученика неверен на строке $abc=001$: $ab+c=1$, а в таблице $y=0$ (единицы только в строках $011,101,110,111$). Верно: $y=ab+ac+bc$. Ход: карта даёт три пары (клетки 3 и 7, 5 и 7, 6 и 7); проверка всегда по строкам с нулём.
  \item \stars{2}~Группа из трёх клеток недопустима (размер не степень двойки). Ответ ученика $F=\overline{A}$ равен единице и на строке $000$, где $F=0$. Верно: клетки 1,3 дают $\overline{A}C$, клетки 2,3 дают $\overline{A}B$, поэтому $F=\overline{A}C+\overline{A}B=\overline{A}(B+C)$.
  \item \stars{3}~$y=d(n+s)$: 2 элемента (ИЛИ и И). Таблица на 8 строк: единица в строках $dn s=101,110,111$. Как проверить: в Конструкторе три входа, переключатели и вкладка Диаграммы; критерий 8 из 8 совпавших строк.
  \item \stars{3}~$y=arm\cdot(ab+ac+bc)$: три И с двумя входами, ИЛИ с тремя и ещё одно И с $arm$ (5 элементов). Вариант без скобок $arm\,ab+arm\,ac+arm\,bc$ --- 4 элемента, но 3-входовых. Как проверить: на вкладке Диаграммы единица ровно в 4 строках из 16 ($arm\,abc=1011,1101,1110,1111$); при $arm=0$ выход нулевой. Образец: урок «Голосование «2 из 3» с охраной».
  \item \stars{3}~При $B=C=1$ первая функция равна $A+\overline{A}$, и из-за задержки инвертора на выходе возникает иголка нуля. Во второй добавлено $BC=1$, которое держит выход в единице весь переход. На карте: группы $AB$ и $\overline{A}C$ только соприкасаются, а $BC$ накрывает обе клетки на границе.
\end{enumerate}

\section*{Минимизация логических функций}
\begin{enumerate}
  \item \stars{1}~$F_1=\overline{B}\,\overline{D}$ (четыре угла); $F_2=D$ (восемь клеток); $F_3=\overline{A}\,\overline{B}$ (верхняя строка). Проверка: соберите $F_1$ в Конструкторе и перебором $0$--$15$ убедитесь, что единицы только в строках 0, 2, 8, 10.
  \item \stars{2}~$F=\overline{D}$: чётные цифры --- это ровно те, у которых $D=0$, а безразличные состояния 10, 12, 14 лишь расширяют группу до восьми клеток.
  \item \stars{2}~Простые импликанты: $\overline{A}\,\overline{B}$ (0,1), $\overline{A}\,\overline{C}$ (0,2), $\overline{B}C$ (1,5), $B\overline{C}$ (2,6), $AC$ (5,7), $AB$ (6,7). Минимальное покрытие состоит из трёх импликант, вариантов два: $\overline{A}\,\overline{B}+B\overline{C}+AC$ и $\overline{A}\,\overline{C}+\overline{B}C+AB$ (циклическое покрытие: существенных импликант нет).
  \item \stars{1}~Например, $F=\overline{B}\,\overline{D}$: четыре угла карты образуют одну группу. Другой пример: $F=\overline{B}\,\overline{C}$ ($\Sigma m(0,1,8,9)$), группа замыкается по вертикали.
  \item \stars{2}~Одна LUT4 (таблица на четыре входа в ПЛИС) реализует любую функцию от четырёх переменных. Число слагаемых на размер не влияет, а функция от пяти и более переменных не поместится в одну LUT.
  \item \stars{1}~Минимальная форма: $F=\overline{A}B+B\overline{C}$ (строка $AB=01$ целиком плюс блок $2\times2$). Букв 4 вместо $6\cdot4=24$ в СДНФ.
  \item \stars{2}~$F=\overline{B}\,\overline{C}+\overline{B}\,\overline{D}+A\overline{C}$. Ход: $\{0,1,8,9\}$ --- $\overline{B}\,\overline{C}$ (замыкание сверху вниз), $\{0,2,8,10\}$ --- четыре угла, $\{8,9,12,13\}$ --- $A\overline{C}$. Клетки 1, 2, 12 и 13 накрыты единственным образом, поэтому все три группы существенны.
  \item \stars{2}~Лишнее слагаемое --- $BD$ (клетки 5, 7, 13, 15): все они уже накрыты четырьмя обязательными парами $\overline{A}\,\overline{C}D$, $\overline{A}BC$, $ACD$, $AB\overline{C}$. Минимальная форма состоит из этих четырёх слагаемых.
  \item \stars{2}~Клетки 3 ($0011$), 6 ($0110$) и 9 ($1001$) попарно не соседние и не образуют прямоугольник. Единица 3 объединяется с безразличной 11 ($\overline{B}CD$), 6 --- с 14 ($BC\overline{D}$), а 9 --- с 11, 13, 15 в квадрат $AD$.
  \item \stars{2}~$F=\overline{A}\,\overline{B}\,\overline{C}\,\overline{D}+\overline{B}CD+BC\overline{D}+AD$. Ход: клетка 0 изолирована (соседи 1, 2, 4, 8 равны нулю), остальные три единицы расширяются безразличными клетками $10$--$15$.
  \item \stars{3}~Обе схемы совпадают на значениях $0$--$9$ (и по условию $10$--$15$ безразличны). По карте: $F=A+B+CD$ --- ИЛИ на три входа и один И (2 блока) против пяти блоков в варианте с компараторами. Как проверить: вкладка Диаграммы, перебор $x=0\ldots15$, сравнение с таблицей $3\le x\le9$ на значениях $0$--$9$; на $10$--$15$ у второй схемы выход $1$, и это допустимо.
  \item \stars{3}~На карте единицы в клетках 11, 13, 14, 15: три группы по две клетки, $y=ABD+ABC+ACD$ (с $A=arm$, $B=a$, $C=b$, $D=c$): 4 элемента, 12 входов; со скобкой $arm\cdot(ab+ac+bc)$ --- 5 элементов, 11 входов. Как проверить: на вкладке Диаграммы единица ровно в 4 строках из 16; оба варианта дают одну таблицу (образец --- урок «Голосование «2 из 3» с охраной»).
  \item \stars{3}~«Хотя бы двое из четырёх»: единицы на всех строках с двумя и более единицами (11 строк из 16). Минимальная форма: $ab+ac+ad+bc+bd+cd$ --- шесть слагаемых по две буквы (6 элементов И и один ИЛИ на шесть входов). Как проверить: в Конструкторе все 16 комбинаций сверяются с картой; единица ровно в 11 строках.
\end{enumerate}

\section*{Дешифраторы и шифраторы}
\begin{enumerate}
  \item \stars{1}~3 входа и 8 выходов; при 101 активна линия с номером 5.
  \item \stars{1}~$y_0=\overline{A_1}\,\overline{A_0}\,EN$, $y_1=\overline{A_1}A_0 EN$, $y_2=A_1\overline{A_0}EN$, $y_3=A_1A_0EN$.
  \item \stars{1}~Горят сегменты $a$, $b$ и $c$.
  \item \stars{1}~Старший активный вход $I_2$, поэтому на выходе 10 (число 2); вход $I_1$ игнорируется.
  \item \stars{1}~$6=110_2$; в коде Грея $101$. В коде Грея соседние числа различаются одним битом, поэтому при переходе не возникает ложных промежуточных значений.
  \item \stars{2}~Пять: один дешифратор первого уровня и четыре второго ($1+4$). Первый выбирает по старшим двум разрядам, какому из четырёх второго уровня дать разрешение $E$.
  \item \stars{2}~Сегмент $c$ не горит только на цифре 2 ($0010$). С безразличной клеткой 10 (вместе с клеткой 2) получаем $\overline{c}=\overline{x_2}\,x_1\overline{x_0}$, откуда $c=x_2+\overline{x_1}+x_0$: два элемента (ИЛИ с тремя входами и инвертор).
  \item \stars{2}~$b_2=g_2=1$, $b_1=b_2\oplus g_1=1\oplus1=0$, $b_0=b_1\oplus g_0=0\oplus0=0$, то есть \code{100} (число 4). Обратно: $g_2=1$, $g_1=1\oplus0=1$, $g_0=0\oplus0=0$ --- \code{110}.
  \item \stars{2}~Верна оценка ученика только для приоритетного шифратора. Схема из двух ИЛИ при $0110$ даёт $y_0=x_1+x_3=1$, $y_1=x_2+x_3=1$, то есть \code{11}. Чтобы получить $10$ (старшая единица $x_2$), нужен приоритетный шифратор: $y_1=x_3+x_2$, $y_0=x_3+\overline{x_2}\,x_1$.
  \item \stars{2}~Код Грея числа 5 ($101$) равен $111$: $g_2=1$, $g_1=1\oplus0=1$, $g_0=0\oplus1=1$. Код $100$ --- код Грея числа 7. «Просто сдвинуть» нельзя: $g=b\oplus(b\gg1)$.
  \item \stars{3}~Цепочка: приоритетный шифратор (ширина 4) $\to$ «Объединение шин» (константа $00$ шириной 2 и номер) $\to$ «BCD$\to$7 сегментов». Коды сегментов: зона 1 --- \code{0000110}, зона 2 --- \code{1011011}, зона 3 --- \code{1001111}, зона 0 --- \code{0111111}; при двух зонах (\code{0110}, \code{1010}) показывается старшая: 2 и 3. Как проверить: в Диаграммах задать шесть значений $z$ и сверить. Для варианта с гашением нужен выход $V$ и элементы И.
  \item \stars{3}~Два дешифратора $2\to4$ получают $a_1a_0$; первому разрешение $\overline{a_2}$, второму $a_2$. Выходы первого --- $y_0\ldots y_3$, второго --- $y_4\ldots y_7$. Как проверить: на вкладке Диаграммы $a$ счётчиком $0$--$7$, единица перебегает по выходам, при $E=0$ все нули. Блоки: «Дешифратор» (с разрешением), «Срез шины», НЕ.
  \item \stars{3}~$g_2=b_2$, $g_1=b_2\oplus b_1$, $g_0=b_1\oplus b_0$ (два XOR). Обратное: $b_2=g_2$, $b_1=b_2\oplus g_1$, $b_0=b_1\oplus g_0$ (два XOR цепочкой). Как проверить: счётчик $0$--$7$ на входе, на каждом шаге меняется ровно один бит кода Грея; на выходе обратного преобразователя исходное число.
\end{enumerate}

\section*{Мультиплексоры и демультиплексоры}
\begin{enumerate}
  \item \textbf{(\stars{1})} 3 адресных входа ($2^3=8$); у $16:1$ их 4 ($2^4=16$). Число адресных входов равно $\log_2$ числа информационных.
  \item \textbf{(\stars{1})} Семь ($4+2+1$): 4 на первом уровне, 2 на втором, 1 на третьем.
  \item \textbf{(\stars{1})} Восемь, по одному на разряд, с общим адресным входом.
  \item \textbf{(\stars{1})} Выбран выход $y_2$: $y_2=1$, остальные три выхода равны 0 (адрес $10_2=2$).
  \item \textbf{(\stars{2})} На $D_0..D_3$ подают $C,\ \overline{C},\ \overline{C},\ C$ (при $AB=00,01,10,11$): при $AB=00$ функция равна $C$, при $01$ и $10$ равна $\overline{C}$, при $11$ равна $C$. Как проверить в Logra: перебор 8 комбинаций на вкладке Диаграммы.
  \item \textbf{(\stars{2})} Любая функция $n$ переменных реализуется мультиплексором с $2^n$ информационными входами, на которые подаются константы из столбца таблицы истинности. Для И подают \code{0001}, для ИЛИ \code{0111} (порядок $AB=00,01,10,11$). \lookahead{так же устроена LUT в ПЛИС.}
  \item \textbf{(\stars{2})} Получилось ИЛИ: комбинация $AB=11$ даёт 1, а нужен 0. Правильные константы $0,1,1,0$.
  \item \textbf{(\stars{2})} Приёмник адресует выходы на один шаг раньше (или позже), чем передатчик: данные каждого канала окажутся на соседнем выходе (газ на «крене», крен на «тангаже» и т.\,д., рыскание на «газе»). Адреса на обоих концах должны быть синхронны.
  \item \textbf{(\stars{3})} Ход: $x_0..x_3$ --- датчики зон, $s_1s_0$ --- номер зоны; демультиплексор с тем же адресом зажигает лампу зоны. Как проверить: на Диаграммах для адреса $00..11$ выход равен соответствующему датчику, а изменение остальных датчиков выход не меняет; лампа номера горит ровно одна.
  \item \textbf{(\stars{3})} Ход: мультиплексор 4:1 с шириной данных 4, на $D_0,D_1,D_2,D_3$ подать $a,b,c,c$. Как проверить: на Диаграммах при $s=0,1,2,3$ выход равен $a,b,c,c$.
  \item \textbf{(\stars{3})} Ход: на $D_0..D_3$ при адресе $AB=00,01,10,11$ подают $0,\ C,\ C,\ 1$ (мажоритарная функция: при $AB=01$ и $10$ результат равен $C$). Как проверить: 8 комбинаций на Диаграммах дают единицу, если единиц на входах не меньше двух.
\end{enumerate}

\section*{Арифметические узлы}
\begin{enumerate}
  \item \textbf{(\stars{1})} $1011_2+0110_2=10001_2$, то есть $11+6=17$.
  \item \textbf{(\stars{1})} $5=00000101_2$, инверсия $11111010$, $+1$ даёт $11111011_2=-5$.
  \item \textbf{(\stars{1})} От $-128$ до $+127$ ($-2^7\ldots 2^7-1$).
  \item \textbf{(\stars{1})} $sum=a\oplus b\oplus c_{in}=1$, $c_{out}=1$ (единиц три, не меньше двух).
  \item \textbf{(\stars{1})} Около $8\times 0{,}5=4$~нс: перенос проходит все разряды по очереди.
  \item \textbf{(\stars{2})} $7-3=7+\overline{3}+1$: $0111+1101=1\,0100$; перенос из старшего разряда отбрасываем, остаётся $0100=4$.
  \item \textbf{(\stars{2})} Единицей станет $A>B$: сравнение идёт от старшего разряда, и в третьем разряде $A=1$, $B=0$ (старшие два совпали).
  \item \textbf{(\stars{2})} $0110+0101=1011_2$; как знаковое это $-5$, а верный ответ $11$ не помещается в диапазон $-8\ldots 7$. $c_3=1$, $c_4=0$, $V=c_3\oplus c_4=1$: переполнение, ответ неверен.
  \item \textbf{(\stars{2})} $\overline{B}=1101$, $1101+1101+1=1\,1011$; результат $1011=-5$, $c_3=1$, $c_4=1$, $V=0$, $N=1$. $N\oplus V=1$: действительно $A<B$ ($-3<2$).
  \item \textbf{(\stars{2})} Забыта единица: $A-B=A+\overline{B}+1$, а студент прибавил только $\overline{B}$. Правильно: $0110+1101+1=1\,0100$, то есть $0100=4$ (на $c_0=1$ вычитатель делает это сам).
  \item \textbf{(\stars{2})} $0111+0001=1000$, как знаковое $-8$ --- неверно. $Z=0$, $N=1$, $C=0$ ($c_4=0$), $V=1$ ($c_3=1$, $c_4=0$): переполнение знаковых чисел, хотя переноса нет.
  \item \textbf{(\stars{3})} Ход: компаратор $A>B$ (блок «Меньше: A < B» с переставленными входами или вычитатель). Как проверить: при $B=8$ лампа горит ровно при $A=9..15$ (вкладка Диаграммы, перебор 16 значений); при $B=15$ ни разу.
  \item \textbf{(\stars{3})} Ход: $s$ = сумма, $c_4$ = перенос; результат = $c_4 ? 1111 : s$, то есть 4-разрядный мультиплексор 2:1 (глава~\ref{ch:mux}). Как проверить: $7+3=10$, $9+9=15$, $15+1=15$.
  \item \textbf{(\stars{3})} Ход: сумматор и вычитатель, элементы И и ИЛИ, мультиплексор 4:1 по коду операции и элемент ИЛИ-НЕ над битами результата (ноль). Как проверить: для $(2,1)$ ожидается $3,1,0,3$; для $(3,3)$ --- $2,0,3,3$ (в 2 битах); для $(0,0)$ все операции дают 0 и $Z=1$.
\end{enumerate}

\section*{Триггеры}
\begin{enumerate}
  \item \stars{1} $S=R=1$: оба выхода становятся нулями, а после снятия сигналов состояние триггера неопределённо. Проверка: урок «SR-триггер: установить и сбросить», шаг «Запрещённая комбинация».
  \item \stars{2} Защёлка прозрачна, пока тактовый сигнал равен 1; триггер запоминает значение только в момент фронта. Ход: при $C=1$ выход защёлки повторяет все изменения $D$, а выход триггера остаётся постоянным до фронта.
  \item \stars{2} $T_{min}=0{,}5+3+0{,}4=3{,}9$~нс, $f_{max}\approx 256$~МГц. Ход: период обязан вместить задержку выхода, логику и время предустановки.
  \item \stars{1} $D=\overline{Q}$ превращает D-триггер в T-триггер; при $T=1$ он делит частоту такта на 2. Проверка: урок «Делитель частоты на 2».
  \item \stars{1} Выход переключается на противоположный.
  \item \stars{2} Метастабильность возникает, когда вход меняется внутри окна setup/hold: триггер «зависает» между 0 и 1. Два триггера подряд дают время успокоиться (но вероятность ошибки не нулевая, и для шины такой приём не годится).
  \item \stars{2} Защёлка при $C=1$ с $D=\overline{Q}$ превращается в кольцо с нечётным числом инверсий, то есть в генератор: $Q$ непрерывно переключается (частота зависит от задержек элементов, а не от такта). Исправление: заменить защёлку D-триггером, который берёт значение по фронту.
  \item \stars{2} Без обратной связи $D$ повторяет датчик: после закрытия двери датчик снова 0, и триггер при следующем фронте запишет 0, то есть тревога не запоминается. Нужно подать на $D$ сигнал «датчик ИЛИ $Q$».
  \item \stars{3} Схема: $D=\text{sensor}\lor Q$, при $\text{key}=1$ синхронный сброс (приоритет у сброса: $D_{\text{new}}=\overline{\text{key}}\land(\text{sensor}\lor Q)$). Как проверить в Logra: Конструктор + Диаграммы; открыть дверь на один такт, убедиться, что $Q$ держится; подать \code{key}; затем подать \code{key} и \code{sensor} вместе: $Q=0$. Образец: проект «Сигнализация с запоминанием».
  \item \stars{3} Цепочка: кнопка $\to$ два триггера (синхронизатор) $\to$ триггер с запоминанием предыдущего значения $\to$ $\text{pulse}=\text{sync}\land\overline{\text{prev}}$. Как проверить в Logra: уроки «Синхронизатор: привязать кнопку к такту» и «Детектор фронта: один импульс на нажатие»; на Диаграммах при кнопке, зажатой на 10 тактов, pulse равен 1 ровно один такт.
\end{enumerate}

\section*{Регистры и счётчики}
\begin{enumerate}
  \item \stars{1} $0010,\ 0100,\ 1000,\ 0001$.
  \item \stars{1} После $1111$ счётчик переходит в $0000$; модуль счёта равен 16.
  \item \stars{2} 25 разрядов ($2^{25}=33\,554\,432\ge 27\,000\,000$). Ход: нужно $\lceil\log_2 27\,000\,000\rceil=25$; для точной секунды считают от 0 до $26\,999\,999$ и по сравнению сбрасывают.
  \item \stars{1} Считает от 0 до 9; после 9 возвращается в 0 по сигналу сравнения $Q=9$.
  \item \stars{2} 15 состояний ($2^4-1$): нулевое состояние в цикл не входит.
  \item \stars{1} Все триггеры тактируются одновременно, поэтому нет цепочки задержек и частота выше; выходы меняются одновременно.
  \item \stars{2} При сбросе по значению 10 ($1010$) счётчик проходит $0,1,\ldots,9,10,0$: 11 состояний, а на мгновение возникает лишнее значение 10. Сбрасывать нужно при $Q=9$ ($N-1$).
  \item \stars{2} Состояние $000$ у LFSR с XOR запретно: из него $q_0=q_1\oplus q_2=0$, регистр остаётся в нуле навсегда. Нужно задать ненулевое начальное значение (сброс в $001$).
  \item \stars{3} Схема: регистр, $+1$, компаратор $Q=N-1$, мультиплексор «0 вместо $Q+1$», дешифратор в семь сегментов. Как проверить в Logra: урок «Считаем от 0 до 5» и «Счётчик на семисегментный индикатор»; на Диаграммах цикл $0..N-1$, значения $N$ нет, перенос равен 1 только при $Q=N-1$.
  \item \stars{3} Ход: при $N=4$ импульс приходит раз в 4 такта (проверяется на Диаграммах), для платы $N=27\,000\,000$, считать до $26\,999\,999$, 25 разрядов. Как проверить в Logra: Диаграммы с малым $N$, затем замена параметра. Для меандра 1~Гц переключать T-триггер импульсом $N=13\,500\,000$.
  \item \stars{3} Ход: счётчик секунд $c$ растёт на импульс «секунда» при $\text{alarm}=1$, насыщается на 9, сбрасывается ключом. Как проверить в Logra: четыре проверки на Диаграммах (рост, остановка на 9, остановка при снятии тревоги, обнуление ключом); образцы: «Счётчик с паузой», «Счётчик со сбросом».
\end{enumerate}

\section*{Синхронные схемы, автоматы и память}
\begin{enumerate}
  \item \stars{1} $T=8+0{,}5+0{,}5=9$~нс, $f_{max}\approx 111$~МГц.
  \item \stars{2} Период сократится примерно вдвое (до $0{,}5+4+0{,}5=5$~нс при двух равных этапах по 4~нс, то есть частота около 200~МГц), но результат появится на один такт позже: задержка в тактах вырастает, пропускная способность (один результат за такт) растёт.
  \item \stars{1} У автомата Мура выходы зависят только от состояния, у автомата Мили ещё и от входов. Детектор \code{101} из главы --- автомат Мура: $y=1$ только в состоянии S3. Чтобы получить $y$ в том же такте, нужен автомат Мили с тремя состояниями S0, S1, S2: $y=1$ на переходе из S2 по $x=1$ (в состояние S1).
  \item \stars{1} Для автомата Мура из главы --- 4 состояния и 2 триггера (двоичное кодирование); для автомата Мили достаточно 3 состояний, тоже 2 триггера.
  \item \stars{1} 8 адресных линий; $256\cdot 8=2048$ бит.
  \item \stars{1} $1\text{К}\cdot 16=16\,384$ бит $=2$~КиБ; 10 адресных и 16 информационных линий.
  \item \stars{2} Ошибка в переходе S3 по $x=1$: верный переход ведёт в S1 (последняя единица может начать новую тройку), а не в S0. На \code{101101} верный автомат выдаёт $y=1$ дважды (после 3-го и после 6-го бита), а ошибочный один раз: после 4-го бита он оказывается в S0 вместо S1, и пара «0, 1» не завершает тройку.
  \item \stars{3} Автомат Мура на 5 состояний (ничего, «1», «11», «110», «1101»). Переходы по $x=0/1$: S0: S0/S1; «1»: S0/«11»; «11»: «110»/«11»; «110»: S0/«1101»; «1101»: S0/«11». Для \code{11011011101} совпадений три (позиции 1--4, 4--7, 8--11), значит $y=1$ три раза. Как проверить в Logra: вкладка Диаграммы, подать поток, сравнить число импульсов $y$ с ожидаемым; перекрытие у позиций 1--4 и 4--7 подтверждает правильность перехода «1101»~--- «11».
  \item \stars{3} Адресный счётчик 3 разряда перебирает слова. Коды сегментов (порядок $gfedcba$, 1 --- сегмент горит): 0 --- \code{0111111} (3F), 1 --- \code{0000110} (06), 2 --- \code{1011011} (5B), 3 --- \code{1001111} (4F), 4 --- \code{1100110} (66), 5 --- \code{1101101} (6D), 6 --- \code{1111101} (7D), 7 --- \code{0000111} (07). Проверка: Диаграммы: за 8 тактов все 8 слов по порядку, замена одного слова меняет ровно один такт периода.
\end{enumerate}

\section*{Проектирование автоматов и очереди FIFO}
\begin{enumerate}
  \item \stars{2} Три состояния (красный, жёлтый, зелёный); при двоичном кодировании 2 триггера, при унитарном 3 (счётчик времени в состоянии не учтён).
  \item \stars{2} Ветка \code{if (c10)} имеет приоритет, пятёрка будет проигнорирована и потеряна. Нужно добавить в таблицу столбец «обе монеты» (в $S_0$ она даёт 15 рублей, то есть сразу продажу).
  \item \stars{2} Состояния: ждём первую 1, увидели 1, увидели 11, увидели 110. Автомат Мура: 4 состояния, у последнего выход 1. Автомат Мили: 3 состояния, выход 1 при входе 0 в состоянии «11».
  \item \stars{1} Указатели 4-разрядные, счётчик слов 5-разрядный (значения 0--16).
  \item \stars{1} Потеряно 3 слова: когда очередь полна, запись игнорируется.
  \item \stars{1} Удобнее: меньше состояний и реакция в том же такте. Опаснее: выход зависит от входа напрямую, поэтому помехи проходят на выход, а входы с выходами образуют длинные комбинационные цепи.
  \item \stars{2} При пустой и при полной очереди указатели равны, поэтому выражения для full и empty совпадут и признаки дадут одно и то же (оба 1 или оба 0). Исправление: хранить счётчик слов (\code{empty}: count=0, \code{full}: count=8) или сделать указатели на один бит шире (пусто, если указатели равны полностью; полна, если младшие биты равны, а старшие различаются).
  \item \stars{2} Сценарий: хозяин открыл дверь ($d=1$) и в тот же такт ввёл верный код ($c=1$). Пустая ячейка оставляет поведение неопределённым: автомат может «зависнуть» или включить тревогу. Нужно записать переход в ОХРАНА\_ВЫКЛ (код приоритетнее датчика), как в таблице главы.
  \item \stars{3} Граф и таблица как в разделе «Сквозной проект»: три состояния, переходы по $(d,c)$, выход «сирена»=1 только в ТРЕВОГА. Проверка на диаграмме: после \code{101} состояние ОХРАНА; $d=1$ в ОХРАНА включает сирену на следующем такте; $d=1$ в ОХРАНА\_ВЫКЛ её не включает; повторный \code{101} в ТРЕВОГА выключает сирену. Вкладка RTL: регистр состояния и логики переходов и выхода; образцы: урок «Кодовый замок» и проект \code{door-controller}.
  \item \stars{3} Счётчик слов \code{used} (0..4); запись разрешена только при $\text{used}<4$, чтение при $\text{used}>0$; $\text{full}=(\text{used}=4)$, $\text{empty}=(\text{used}=0)$. При пяти записях подряд указатель записи останавливается на 0 (после 4 записей обернулся), \code{used}=4, пятое слово игнорируется, первое слово не затёрто. При пяти чтениях подряд пятое ничего не меняет, \code{used}=0. Проверка: вкладка Диаграммы (пример \code{fifo-four}).
\end{enumerate}

\section*{Электрические основы: КМОП, выходы и шины}
\begin{enumerate}
  \item \stars{2} Три pMOS последовательно между питанием и выходом и три nMOS параллельно между выходом и землёй: 6 транзисторов (3 pMOS и 3 nMOS).
  \item \stars{1} Мощность уменьшится в $(3{,}3/1{,}8)^2\approx 3{,}4$ раза.
  \item \stars{1} $P = 10\cdot 10^{-12}\cdot 3{,}3^2\cdot 27\cdot 10^6 \approx 2{,}9$~мВт.
  \item \stars{1} $N_H = 2{,}6 - 2{,}0 = 0{,}6$~В.
  \item \stars{2} Открытый сток исключает конфликты: любое устройство может только опустить линию в 0, а единицу даёт резистор. Получается монтажное И, и ведомый может задержать ведущего.
  \item \stars{1} Один буфер тянет линию к питанию, другой к земле: через них течёт сквозной ток, а напряжение на линии неопределённо (в симуляции \code{X}); возможен перегрев.
  \item \stars{2} Два двухтактных выхода, соединённые напрямую, при разных значениях дают сквозной ток (один тянет к питанию, другой к земле) и неопределённый уровень. Исправления: (1) завести оба сигнала на логический элемент ИЛИ; (2) использовать выходы с открытым стоком и общей подтяжкой (получится монтажное И инверсных сигналов) или буферы с тремя состояниями, у которых разрешён только один.
  \item \stars{2} $R=(3{,}3-2{,}0)/0{,}008=162{,}5$~Ом. Ближайший меньший номинал 150~Ом дал бы ток $1{,}3/150\approx 8{,}7$~мА, больше допустимого, поэтому берём 180~Ом: $I=1{,}3/180\approx 7{,}2$~мА. Как проверить в Logra: урок «Светодиод и ограничение тока» и вкладка Лаборатория.
  \item \stars{3} Линия равна 1, если не нажата ни одна кнопка, и 0, если нажата хотя бы одна: таблица (не нажаты, нажата A, нажата B, обе) даёт $1,0,0,0$, то есть монтажное И уровней (И-НЕ нажатий). Ток через подтяжку при нажатых обеих кнопках $3{,}3/10\,000=0{,}33$~мА, так как резистор один и ток от числа нажатых кнопок не зависит. Проверка: вкладки Выводы (подтяжка «к 1») и Диаграммы.
  \item \stars{3} Из времени нарастания $2{,}2\cdot R\cdot 100\cdot 10^{-12}\le 10^{-6}$ получаем $R\le 4{,}55$~кОм. Из тока $3{,}3/R\le 0{,}8\cdot 10^{-3}$ получаем $R\ge 4{,}1$~кОм. Подходит диапазон 4,1--4,55~кОм, например 4,3~кОм из ряда E24: время нарастания $2{,}2\cdot4300\cdot10^{-10}\approx 0{,}95$~мкс, ток $\approx 0{,}77$~мА. Оба условия выполняются, но с малым запасом.
\end{enumerate}

\section*{Цифра и аналоговый мир}
\begin{enumerate}
  \item \stars{1} $3{,}3/1024\approx 3{,}2$~мВ.
  \item \stars{1} $5/0{,}01 = 500$ уровней, нужно 9 бит ($2^9=512$).
  \item \stars{1} Не меньше 6~кГц (вдвое больше наивысшей частоты).
  \item \stars{1} $192/256=75\,\%$; среднее напряжение $0{,}75\cdot 3{,}3\approx 2{,}5$~В.
  \item \stars{2} 12 тактов; параллельному АЦП нужно $2^{12}-1 = 4095$ компараторов.
  \item \stars{2} Пробные коды ЦАП: \code{0x80}, \code{0x40}, \code{0x60}, \code{0x70}, \code{0x68}, \code{0x64}, \code{0x66}, \code{0x65}. Результат \code{0x64}, то есть 100.
  \item \stars{2} При ответе «нет» (вход меньше пробного кода) бит нужно \emph{сбросить}: результат остаётся \code{0x00}, а не \code{0x80}. Вторая пробная комбинация \code{0x40}.
  \item \stars{3} Восьмиразрядный счётчик считает $0,1,\ldots,255$ и обнуляется; выход $\text{pwm}=(\text{counter}<\text{level})$. При $\text{level}=64$ единица в течение 64 тактов из 256 (25\,\%), при $\text{level}=0$ условие никогда не истинно, при $\text{level}=255$ выход равен 0 только на последнем такте. Проверка: вкладка Диаграммы (пример \code{pwm}).
  \item \stars{3} Период $27\,000\,000/2000=13\,500$ тактов; единица $0{,}25\cdot 13\,500=3375$ тактов, ноль 10\,125. Схема: счётчик по модулю 13\,500 (14 разрядов), выход $=(\text{counter}<3375)\ \&\ \text{тревога}$. Проверка: Диаграммы: период 13\,500 тактов, после выхода из тревоги выход равен нулю; затем зуммер на вкладке Подключение.
\end{enumerate}

\section*{Тактирование и временной анализ}
\begin{enumerate}
  \item \stars{1}~$T = 0{,}5+4+3+0{,}5 = 8$~нс, $f_{max}=125$~МГц.
  \item \stars{1}~Запас $10-8{,}2=1{,}8$~нс. При 150~МГц период 6,67~нс и запас отрицательный ($-1{,}5$~нс): схема не уложится.
  \item \stars{2}~Считать до $27\,000-1 = 26\,999$.
  \item \stars{1}~$27\cdot 4/1 = 108$~МГц.
  \item \stars{2}~Биты могут прийти с разной задержкой, и приёмник на мгновение увидит значение, которого не было (при переходе \code{0111}$\to$\code{1000} возможны \code{1111} или \code{0000}). Нужны код Грея, квитирование или FIFO.
  \item \stars{2}~$N = 0{,}005\cdot 27\cdot 10^6 = 135\,000$ тактов.
  \item \stars{2}~Период при 100~МГц равен 10~нс, путь 12~нс: запас $10-12=-2$~нс, схема не уложится. Предельная частота около $1/12~\text{нс}\approx 83$~МГц. Исправить: разбить логику конвейером, поставить PLL на 80~МГц или меньше. Как проверить: «Кристалл» с частотой 100~МГц покажет невыполнение требования.
  \item \stars{2}~Выход триггера идёт по обычным проводам, а не по тактовой сети: у разных триггеров он приходит с разной задержкой (перекос), и в симуляции без задержек этого не видно. Правильно: \code{reg en; always @(posedge clk) en <= ~en; always @(posedge clk) if (en) ...;} --- один \code{clk}, а медленный блок разрешается импульсом.
  \item \stars{3}~Проект. Как проверить: на «Кристалле» $f_{max}\ge 54$~МГц; в диаграмме серия \code{1-0-1-0-1} на кнопке после синхронизатора и фильтра даёт один импульс \code{code\_step}; \code{rst} приводит \code{state} в OFF на ближайшем фронте.
  \item \stars{3}~Проект. Делитель до $13\,500\,000-1$ даёт \code{tick} раз в 0,5~с; \code{if (tick) led <= ~led;}. Как проверить: на RTL один \code{clk} на всех триггерах; на диаграмме (с уменьшенной константой) \code{led} меняется только при \code{tick = 1}; на «Кристалле» $f_{max}$ не ниже 100~МГц.
\end{enumerate}

\section*{Последовательные интерфейсы}
\begin{enumerate}
  \item \stars{1}~$27\,000\,000/9600 \approx 2812{,}5$, то есть около 2813 тактов.
  \item \stars{1}~$10\cdot 8{,}68 = 86{,}8$~мкс; $115\,200/10 = 11\,520$ байт в секунду.
  \item \stars{1}~Старт 0, затем биты $1,0,1,0,1,0,1,0$ (младший первым), стоп 1: линия строго чередуется \code{0101010101}.
  \item \stars{2}~В середине сигнал устоялся и дальше всего от фронтов, поэтому небольшие расхождения скоростей и помехи не приведут к ошибке.
  \item \stars{2}~Да: за кадр накопится ошибка $9{,}5\cdot 3\,\% \approx 0{,}29$ бита, меньше половины. Допустимое расхождение скоростей около $0{,}5/9{,}5\approx 5\,\%$.
  \item \stars{2}~У SPI есть общий такт, нет стартовых и стоповых бит, и данные идут в обе стороны одновременно. Недостаток: больше проводов (по линии CS на каждое ведомое устройство) и нет адресации и подтверждения.
  \item \stars{2}~Ошибка: выборка идёт по границам битов, там сигнал меняется. Любая задержка фронта или помеха на длинном проводе даёт чтение соседнего бита. Исправление: после спада линии отсчитать полбита (середина стартового бита), проверить, что линия всё ещё 0, и дальше читать каждые \code{CLKS\_PER\_BIT} тактов. Как проверить в Logra: «UART-приёмник 8N1», кадр \code{0x41} с небольшим сдвигом фронта.
  \item \stars{3}~Проект. Раз в секунду (счётчик до $27\,000\,000-1$) подаётся импульс \code{start}; в \code{uart\_tx} данные \code{state} (0, 1 или 2), \code{CLKS\_PER\_BIT=234}. Как проверить: на диаграмме (с \code{CLKS\_PER\_BIT=8}) стартовый 0, 8 бит младшим вперёд, стоповая 1; на плате в «UART-терминале» вкладки «Подключение» раз в секунду приходит символ.
  \item \stars{3}~Проект. Выход \code{tx} передатчика подаётся на \code{rx} приёмника, \code{rx\_data} выводится на светодиоды. Как проверить: перебор 256 значений на диаграмме, у каждого \code{rx\_data} совпадает с отправленным; при стоповом бите 0 сигнал \code{rx\_valid} не поднимается (кадр отбрасывается).
\end{enumerate}

\section*{Память SDRAM на 64 Мбит}
\begin{enumerate}
  \item \stars{1}~$4\cdot 2048\cdot 256 = 2\,097\,152=2^{21}$ слов, адрес из 21 бита (2 бита банка, 11 строки, 8 столбца).
  \item \stars{1}~$133\cdot 10^6\cdot 4~\text{байта}\approx 532$~Мбайт/с.
  \item \stars{2}~$64~\text{мс}/2048=31{,}25$~мкс; при 100~МГц это 3125 тактов.
  \item \stars{2}~Строка ещё не открыта, пока не прошло $t_\text{RCD}$ тактов: данные будут неверными или неопределёнными. Автомат должен после ACTIVE подождать два такта (NOP) и только затем выдать READ.
  \item \stars{3}~Проект. \code{assign \{ba, row, col\} = addr;} и \code{case (cmd)} по таблице команд; как проверить: для \code{0x1FFFFF} $ba=3$, $row=0x7FF$, $col=0xFF$; на диаграмме все семь команд дают комбинации таблицы, ветка \code{default} выдаёт NOP.
  \item \stars{3}~Проект. Счётчик адреса 21 бит, регенерация раз в 31,25~мкс (3125 тактов при 100~МГц), запись по событию. Время заполнения: $2^{21}$ слов; при одном событии в секунду около 24 суток. Как проверить: на модели с малыми константами счётчик возвращается в 0 после максимального адреса.
\end{enumerate}

\section*{Собираем процессор}
\begin{enumerate}
  \item \stars{1}~Цикл занимает 9 команд по два такта: 18 тактов.
  \item \stars{2}~\code{LDI 1; STA 0; LDI 5; OUT; SUB 0; JZ 8; JMP 3; NOP; OUT; HLT}, то есть коды \code{11 30 15 80 50 78 63 00 80 90}. Константу 1 нужно сначала записать в ОЗУ, потому что при включении память не обнулена. На выходе появятся 5, 4, 3, 2, 1, 0.
  \item \stars{1}~16 ячеек. Чтобы увеличить программу вдвое, нужен 5-разрядный счётчик команд и более широкая команда (адрес перехода должен иметь 5 бит).
  \item \stars{2}~Последнее число, помещающееся в 8 бит, --- 233 (четырнадцатое по счёту, $F_{13}$); следующее $377 - 256 = 121$ появится после переполнения аккумулятора.
  \item \stars{1}~В гарвардской архитектуре память программ и память данных раздельны (свои шины), в архитектуре фон Неймана память общая.
  \item \stars{2}~Читать ПЗУ комбинационно и выполнять команду в том же такте (одноцикловый процессор); критический путь при этом станет длиннее, а частота ниже.
  \item \stars{2}~Команда \code{JMP 2} возвращает счётчик на \code{LDI 5}: аккумулятор каждый раз снова становится 5, \code{SUB} даёт 4, \code{JZ} не срабатывает, на \code{OUT} снова выводится 5. Как проверить: на диаграмме \code{pc} крутится по кругу 2--6, \code{out} всегда 5.
  \item \stars{3}~В \code{case (op)} добавить \code{4'd10: if (acc != 8'd0) pc <= arg;}. Программа \code{15 30 11 31 40 80 51 30 A4 90} (\code{LDI 5; STA 0; LDI 1; STA 1; LDA 0; OUT; SUB 1; STA 0; JNZ 4; HLT}). Как проверить: на диаграмме \code{out} принимает 5, 4, 3, 2, 1, затем \code{pc} застывает на адресе 9.
  \item \stars{3}~Проект. Программа \code{13 30 11 31 40 80 51 30 A4 80 90} (\code{LDI 3; STA 0; LDI 1; STA 1; LDA 0; OUT; SUB 1; STA 0; JNZ 4; OUT; HLT}): \code{out} принимает 3, 2, 1 и в конце 0, после чего \code{pc} перестаёт меняться. Как проверить: на диаграмме после третьей «попытки» \code{out = 0} и \code{halted = 1}.
\end{enumerate}

